\documentclass[11pt]{article}

\usepackage[a4paper,margin=1in]{geometry}
\usepackage{amsmath,amssymb,amsthm,mathtools}
\usepackage{enumitem}

\newtheorem{theorem}{Theorem}[section]
\newtheorem{lemma}[theorem]{Lemma}
\newtheorem{proposition}[theorem]{Proposition}
\newtheorem{remark}[theorem]{Remark}
\newtheorem{corollary}[theorem]{Corollary}

\newcommand{\N}{\mathbb N}
\newcommand{\Pow}{\mathcal P}
\newcommand{\HJ}{\operatorname{HJ}}
\newcommand{\FU}{\operatorname{FU}}
\newcommand{\FS}{\operatorname{FS}}

\title{The Graham--Rothschild Theorem for Combinatorial Lines\\
from the Hales--Jewett Theorem}
\author{}
\date{}

\begin{document}

\maketitle

\section{Introduction}

Let \(A\) be a finite alphabet with \(|A|=k\), and let \(v\) be a symbol
not belonging to \(A\). A \emph{variable word} over \(A\) is a word over
\(A\cup\{v\}\) in which \(v\) occurs at least once. If
\[
  w(v)\in (A\cup\{v\})^N
\]
is a variable word and \(a\in A\), we write \(w(a)\in A^N\) for the word
obtained by replacing every occurrence of \(v\) by \(a\). The set
\[
  L(w)=\{w(a):a\in A\}
\]
is called the \emph{combinatorial line} generated by \(w\).

More generally, let \(v_1,\dots,v_m\) be distinct symbols outside \(A\).
An \(m\)-variable word is a word
\[
  z(v_1,\dots,v_m)\in
  \bigl(A\cup\{v_1,\dots,v_m\}\bigr)^N
\]
in which every \(v_i\) occurs at least once. It generates the
\(m\)-dimensional combinatorial subspace
\[
  V(z)
  =
  \left\{
    z(a_1,\dots,a_m):
    (a_1,\dots,a_m)\in A^m
  \right\}.
\]

A combinatorial line contained in \(V(z)\) is obtained by choosing a
variable word
\[
  (b_1,\dots,b_m)\in (A\cup\{v\})^m
\]
and considering
\[
  \left\{
    z\bigl(b_1(a),\dots,b_m(a)\bigr):a\in A
  \right\}.
\]
Equivalently, the line is represented by the variable word
\[
  z(b_1,\dots,b_m),
\]
where at least one of the \(b_i\)'s is equal to \(v\).

We prove the following \(n=1\) instance of the Graham--Rothschild
parameter-sets theorem.

\begin{theorem}[Graham--Rothschild theorem for lines]
  \label{thm:GR-lines}
  For every \(k,m,r\in\N\), with \(k\geq 2\), there exists an integer
  \[
    \operatorname{GR}_1(k,m,r)
  \]
  such that the following holds.

  Let \(A\) be an alphabet of size \(k\), let
  \(N\geq \operatorname{GR}_1(k,m,r)\), and let the combinatorial lines
  of \(A^N\) be colored with \(r\) colors. Then there exists an
  \(m\)-dimensional combinatorial subspace \(V\subseteq A^N\) such that
  all combinatorial lines contained in \(V\) have the same color.
\end{theorem}

The proof uses the ordinary finite Hales--Jewett theorem, together with
elementary finite Ramsey theory. The argument has two ingredients.

\begin{enumerate}[label=\arabic*.]
  \item Repeated applications of the Hales--Jewett theorem produce a
    structured family in which the color of a variable word depends
    only on the set of blocks in which the variable occurs.
  \item A finite-unions theorem is then used to homogenize all possible
    nonempty unions of those sets of blocks.
\end{enumerate}

\section{A finite-unions theorem}

For positive integers \(x_1,\dots,x_q\), define
\[
  \FS(x_1,\dots,x_q)
  =
  \left\{
    \sum_{i\in F}x_i:
    \varnothing\neq F\subseteq [q]
  \right\}.
\]

For pairwise disjoint nonempty finite sets
\(\gamma_1,\dots,\gamma_m\), define
\[
  \FU(\gamma_1,\dots,\gamma_m)
  =
  \left\{
    \bigcup_{j\in J}\gamma_j:
    \varnothing\neq J\subseteq[m]
  \right\}.
\]

We write
\[
  \gamma_1<\gamma_2<\cdots<\gamma_m
\]
if
\[
  \max\gamma_i<\min\gamma_{i+1}
  \qquad
  \text{for every }i<m.
\]

\subsection{Hales--Jewett implies van der Waerden}

We first record a standard consequence of the Hales--Jewett theorem.

\begin{lemma}
  \label{lem:vdW}
  For every \(\ell,r\in\N\), there exists \(W(\ell,r)\) such that every
  \(r\)-coloring of \([W(\ell,r)]\) contains a monochromatic arithmetic
  progression of length \(\ell\).
\end{lemma}

\begin{proof}
  Let
  \[
    H=\HJ(\ell,r),
  \]
  where the Hales--Jewett theorem is applied to an alphabet of size
  \(\ell\), and put
  \[
    W(\ell,r)=\ell^H.
  \]

  Suppose that
  \[
    c:[W(\ell,r)]\longrightarrow[r]
  \]
  is an \(r\)-coloring. Color a word
  \[
    (a_0,\dots,a_{H-1})\in\{0,\dots,\ell-1\}^H
  \]
  by
  \[
    \widetilde c(a_0,\dots,a_{H-1})
    =
    c\left(
      1+\sum_{i=0}^{H-1}a_i\ell^i
    \right).
  \]

  By the Hales--Jewett theorem, there is a monochromatic variable word.
  Let \(I\subseteq\{0,\dots,H-1\}\) be its nonempty set of variable
  coordinates. Replacing the variable by \(t\in\{0,\dots,\ell-1\}\)
  produces integers of the form
  \[
    a+td,
    \qquad
    0\leq t\leq \ell-1,
  \]
  where
  \[
    d=\sum_{i\in I}\ell^i>0.
  \]
  Thus
  \[
    a,a+d,\dots,a+(\ell-1)d
  \]
  is a monochromatic arithmetic progression of length \(\ell\).
\end{proof}

\subsection{A finite-sums theorem}

We next derive the finite version of Hindman's finite-sums theorem
that will be needed below.

\begin{lemma}
  \label{lem:max-dependent}
  For every \(q,r\in\N\), there exists \(M(q,r)\) such that every
  \(r\)-coloring
  \[
    c:[M(q,r)]\longrightarrow[r]
  \]
  admits positive integers \(x_1,\dots,x_q\) satisfying
  \[
    \sum_{i=1}^q x_i\leq M(q,r)
  \]
  and such that
  \[
    c\left(\sum_{i\in F}x_i\right)
  \]
  depends only on \(\max F\), for every nonempty \(F\subseteq[q]\).
\end{lemma}

\begin{proof}
  We argue by induction on \(q\).

  For \(q=1\), the statement is immediate.

  Assume it has been proved for \(q\). Let
  \[
    M_q=M(q,r).
  \]
  By Lemma~\ref{lem:vdW}, choose \(M(q+1,r)\) sufficiently large that
  every \(r\)-coloring of \([M(q+1,r)]\) contains a monochromatic
  arithmetic progression of length \(M_q+1\), say
  \[
    a,a+d,\dots,a+M_qd.
  \]

  Let
  \[
    c:[M(q+1,r)]\longrightarrow[r]
  \]
  be a coloring, and choose such a monochromatic progression. Define
  \[
    c'(s)=c(ds),
    \qquad s\in[M_q].
  \]
  By the induction hypothesis, there exist positive integers
  \(y_1,\dots,y_q\) such that
  \[
    \sum_{i=1}^q y_i\leq M_q
  \]
  and
  \[
    c'\left(\sum_{i\in F}y_i\right)
  \]
  depends only on \(\max F\).

  Set
  \[
    x_i=dy_i
    \quad (1\leq i\leq q),
    \qquad
    x_{q+1}=a.
  \]

  If \(F\subseteq[q]\) is nonempty, then
  \[
    c\left(\sum_{i\in F}x_i\right)
    =
    c'\left(\sum_{i\in F}y_i\right),
  \]
  so its color depends only on \(\max F\).

  If \(q+1\in F\), then
  \[
    \sum_{i\in F}x_i
    =
    a+d\sum_{i\in F\setminus\{q+1\}}y_i.
  \]
  Since
  \[
    0\leq
    \sum_{i\in F\setminus\{q+1\}}y_i
    \leq M_q,
  \]
  this sum belongs to the chosen monochromatic arithmetic progression.
  Consequently all sums whose largest index is \(q+1\) have the same
  color. This completes the induction.
\end{proof}

\begin{proposition}[Finite-sums theorem]
  \label{prop:finite-sums}
  For every \(m,r\in\N\), there exists \(S(m,r)\) such that every
  \(r\)-coloring
  \[
    c:[S(m,r)]\longrightarrow[r]
  \]
  admits positive integers \(z_1,\dots,z_m\) such that
  \[
    \FS(z_1,\dots,z_m)
  \]
  is monochromatic.
\end{proposition}

\begin{proof}
  Let
  \[
    q=(m-1)r+1
  \]
  and let \(S(m,r)=M(q,r)\), where \(M(q,r)\) is given by
  Lemma~\ref{lem:max-dependent}.

  Apply that lemma to obtain \(x_1,\dots,x_q\). Color each index
  \(i\in[q]\) by the color \(c(x_i)\). Since \(q=(m-1)r+1\), the
  pigeonhole principle gives indices
  \[
    i_1<\cdots<i_m
  \]
  for which
  \[
    c(x_{i_1})=\cdots=c(x_{i_m}).
  \]
  Set
  \[
    z_j=x_{i_j},
    \qquad 1\leq j\leq m.
  \]

  Let \(J\subseteq[m]\) be nonempty. By
  Lemma~\ref{lem:max-dependent}, the color of
  \[
    \sum_{j\in J}z_j
    =
    \sum_{j\in J}x_{i_j}
  \]
  depends only on the largest index \(i_{\max J}\). This is the same as
  the color of \(x_{i_{\max J}}\). Since all the \(x_{i_j}\)'s have the
  same color, every nonempty subset sum of \(z_1,\dots,z_m\) has that
  color.
\end{proof}

\subsection{The finite-unions lemma}

\begin{lemma}[Finite-unions lemma]
  \label{lem:finite-unions}
  For every \(m,r\in\N\), there exists \(L=L(m,r)\) such that every
  coloring
  \[
    d:\Pow([L])\setminus\{\varnothing\}\longrightarrow[r]
  \]
  admits nonempty sets
  \[
    \gamma_1<\gamma_2<\cdots<\gamma_m
  \]
  for which
  \[
    \FU(\gamma_1,\dots,\gamma_m)
  \]
  is monochromatic.
\end{lemma}

\begin{proof}
  Let
  \[
    S=S(m,r)
  \]
  be the number supplied by Proposition~\ref{prop:finite-sums}.

  By repeated applications of the finite Ramsey theorem, choose \(L\)
  large enough that every \(r\)-coloring of the nonempty subsets of
  \([L]\) admits an \(S\)-element set
  \[
    H=\{h_1<\cdots<h_S\}
  \]
  with the following property: for every \(s\in[S]\), all \(s\)-element
  subsets of \(H\) have the same color.

  Let
  \[
    d:\Pow([L])\setminus\{\varnothing\}\longrightarrow[r]
  \]
  be a coloring, and choose such a set \(H\). For \(s\in[S]\), let
  \(e(s)\) denote the common color of all \(s\)-element subsets of \(H\).
  Thus
  \[
    e:[S]\longrightarrow[r]
  \]
  is an \(r\)-coloring.

  Apply Proposition~\ref{prop:finite-sums} to \(e\). We obtain positive
  integers \(x_1,\dots,x_m\) such that
  \[
    x_1+\cdots+x_m\leq S
  \]
  and every nonempty subset sum of the \(x_i\)'s has the same \(e\)-color.

  Partition the first \(x_1+\cdots+x_m\) elements of \(H\) into
  successive blocks
  \[
    \gamma_1<\gamma_2<\cdots<\gamma_m
  \]
  with
  \[
    |\gamma_j|=x_j.
  \]
  For every nonempty \(J\subseteq[m]\),
  \[
    \left|\bigcup_{j\in J}\gamma_j\right|
    =
    \sum_{j\in J}x_j.
  \]
  Since the union is a subset of \(H\), its \(d\)-color is
  \[
    e\left(\sum_{j\in J}x_j\right).
  \]
  The latter is independent of the nonempty set \(J\). Hence all
  members of
  \[
    \FU(\gamma_1,\dots,\gamma_m)
  \]
  have the same color.
\end{proof}

\section{Making the coloring insensitive to constant letters}

The next lemma is the main iterated Hales--Jewett step.

\begin{lemma}[Location-insensitivity lemma]
  \label{lem:location-insensitivity}
  Let \(A\) be a finite alphabet, and let \(r,L\in\N\). There exist
  positive integers
  \[
    M_1,\dots,M_L
  \]
  with the following property.

  Set
  \[
    M=M_1+\cdots+M_L.
  \]
  For every \(r\)-coloring
  \[
    \chi:
    \left\{
      w\in(A\cup\{v\})^M:
      v\text{ occurs in }w
    \right\}
    \longrightarrow[r],
  \]
  there exist variable words
  \[
    w_i(x)\in(A\cup\{x\})^{M_i},
    \qquad i=1,\dots,L,
  \]
  such that the color
  \[
    \chi\bigl(w_1(y_1)\cdots w_L(y_L)\bigr),
    \qquad
    (y_1,\dots,y_L)\in(A\cup\{v\})^L,
  \]
  depends only on the nonempty set
  \[
    I(y_1,\dots,y_L)
    =
    \{i\in[L]:y_i=v\}.
  \]
\end{lemma}

\begin{proof}
  The proof is a backward induction using the Hales--Jewett theorem.

  We choose \(M_1,\dots,M_L\) large enough so that the argument below
  can be carried out. It is enough, for example, to choose each \(M_i\)
  to be a Hales--Jewett number for an alphabet of size \(|A|\) and a
  number of colors large enough to encode every coloring profile that
  can arise from the coordinates before and after the \(i\)-th block.

  For clarity, we describe the backward construction explicitly.

  Suppose that the blocks after the \(i\)-th block have already been
  replaced by single meta-coordinates. At the \(i\)-th stage, fix all
  possible contexts
  \[
    (p,q),
  \]
  where \(p\) is a word in the preceding coordinates and \(q\) is a word
  in the following meta-coordinates, subject to the condition that the
  concatenation \(pq\) already contains the variable symbol \(v\).

  For a constant word \(u\in A^{M_i}\), define its color profile by
  \[
    \mathcal C_i(u)
    =
    \bigl(
      \chi_i(puq)
    \bigr)_{(p,q)},
  \]
  where \(\chi_i\) denotes the coloring present at this stage. There are
  only finitely many possible contexts, and hence only finitely many
  possible profiles.

  Choose \(M_i\) large enough that the Hales--Jewett theorem applies to
  the profile-coloring
  \[
    u\longmapsto\mathcal C_i(u)
  \]
  of \(A^{M_i}\). We obtain a variable word
  \[
    w_i(x)\in(A\cup\{x\})^{M_i}
  \]
  such that
  \[
    \mathcal C_i(w_i(a))
  \]
  is independent of \(a\in A\). Equivalently,
  \[
    \chi_i\bigl(pw_i(a)q\bigr)
  \]
  is independent of \(a\in A\) whenever the context \(pq\) already
  contains \(v\).

  Replace the \(i\)-th block by the meta-coordinate \(x\), and continue
  backward for
  \[
    i=L,L-1,\dots,1.
  \]
  At the end, we obtain variable words
  \[
    w_1(x),\dots,w_L(x)
  \]
  such that changing a constant meta-coordinate from one letter of \(A\)
  to another does not change the color, provided some other
  meta-coordinate is equal to \(v\).

  Now let
  \[
    y=(y_1,\dots,y_L),
    \qquad
    y'=(y_1',\dots,y_L')
  \]
  be two variable meta-words with
  \[
    \{i:y_i=v\}=\{i:y_i'=v\}\neq\varnothing.
  \]
  The words \(y\) and \(y'\) differ only at coordinates at which both
  entries belong to \(A\). We may pass from \(y\) to \(y'\) by changing
  one such constant coordinate at a time. Throughout this process, at
  least one coordinate remains equal to \(v\). By the construction
  above, none of these changes affects the color. Therefore
  \[
    \chi\bigl(w_1(y_1)\cdots w_L(y_L)\bigr)
    =
    \chi\bigl(w_1(y_1')\cdots w_L(y_L')\bigr).
  \]
  Thus the color depends only on the set of variable meta-coordinates.
\end{proof}

\begin{remark}
  The restriction in the construction to contexts already containing
  \(v\) is essential. It guarantees that the intermediate words to
  which the coloring is applied are variable words and therefore lie in
  the domain of \(\chi\).
\end{remark}

\section{Proof of the Graham--Rothschild theorem for lines}

\begin{proof}[Proof of Theorem~\ref{thm:GR-lines}]
  Fix \(k,m,r\in\N\), and let \(A\) be an alphabet of size \(k\).

  Choose \(L=L(m,r)\) as in the finite-unions lemma,
  Lemma~\ref{lem:finite-unions}. Next apply the
  location-insensitivity lemma, Lemma~\ref{lem:location-insensitivity},
  with this value of \(L\). Let
  \[
    M=M_1+\cdots+M_L
  \]
  be the resulting word length.

  We prove that \(M\) has the required Graham--Rothschild property.

  Let the combinatorial lines of \(A^M\) be colored with \(r\) colors.
  Equivalently, color each variable word
  \[
    u(v)\in(A\cup\{v\})^M
  \]
  by the color of the line
  \[
    L(u)=\{u(a):a\in A\}.
  \]
  Denote this coloring by \(\chi\).

  By Lemma~\ref{lem:location-insensitivity}, there exist variable words
  \[
    w_i(x)\in(A\cup\{x\})^{M_i},
    \qquad i=1,\dots,L,
  \]
  such that the color of
  \[
    w_1(y_1)\cdots w_L(y_L)
  \]
  depends only on the nonempty set of indices \(i\) for which \(y_i=v\).

  Fix an arbitrary letter \(a_0\in A\). Define a coloring
  \[
    d:\Pow([L])\setminus\{\varnothing\}\longrightarrow[r]
  \]
  by
  \[
    d(I)
    =
    \chi\bigl(
      w_1(y_1)\cdots w_L(y_L)
    \bigr),
  \]
  where
  \[
    y_i=
    \begin{cases}
      v,& i\in I,\\
      a_0,& i\notin I.
    \end{cases}
  \]
  The location-insensitivity property implies that \(d(I)\) is the
  color of every meta-word whose set of variable coordinates is exactly
  \(I\).

  Apply Lemma~\ref{lem:finite-unions} to \(d\). We obtain nonempty sets
  \[
    \gamma_1<\gamma_2<\cdots<\gamma_m
  \]
  such that
  \[
    d\left(\bigcup_{j\in J}\gamma_j\right)
  \]
  has one fixed value, say \(c_0\), for every nonempty
  \(J\subseteq[m]\).

  Define an \(m\)-variable word \(z(v_1,\dots,v_m)\) of length \(M\) by
  \[
    z(v_1,\dots,v_m)
    =
    w_1(t_1)\cdots w_L(t_L),
  \]
  where
  \[
    t_i=
    \begin{cases}
      v_j,& i\in\gamma_j,\\
      a_0,&
      i\notin\gamma_1\cup\cdots\cup\gamma_m.
    \end{cases}
  \]
  This definition is unambiguous because the \(\gamma_j\)'s are
  pairwise disjoint. Moreover, every variable \(v_j\) occurs, since
  \(\gamma_j\neq\varnothing\). Thus \(z\) generates an
  \(m\)-dimensional combinatorial subspace
  \[
    V
    =
    \left\{
      z(a_1,\dots,a_m):
      (a_1,\dots,a_m)\in A^m
    \right\}
    \subseteq A^M.
  \]

  We claim that every combinatorial line contained in \(V\) has color
  \(c_0\).

  Indeed, a combinatorial line in \(V\) is represented by a variable
  meta-word
  \[
    (b_1,\dots,b_m)\in(A\cup\{v\})^m,
  \]
  where at least one \(b_j\) is equal to \(v\). Its corresponding
  variable word in \(A^M\) is
  \[
    z(b_1,\dots,b_m).
  \]
  At the level of the \(L\) blocks, the variable \(v\) occurs precisely
  in the blocks indexed by
  \[
    I_b
    =
    \bigcup_{\{j\in[m]:b_j=v\}}\gamma_j.
  \]
  This is a nonempty union of the \(\gamma_j\)'s. By the
  location-insensitivity property,
  \[
    \chi\bigl(z(b_1,\dots,b_m)\bigr)=d(I_b).
  \]
  By the choice of the \(\gamma_j\)'s,
  \[
    d(I_b)=c_0.
  \]
  Therefore every combinatorial line contained in \(V\) has color
  \(c_0\).

  This proves the theorem for words of length \(M\). If \(N>M\), fix
  \(N-M\) coordinates to the letter \(a_0\) and apply the length-\(M\)
  case to the remaining coordinates. Hence the conclusion holds for
  every \(N\geq M\).

  We may therefore take
  \[
    \operatorname{GR}_1(k,m,r)=M.
  \]
\end{proof}

\section{The two-color formulation}

The form commonly used in density Hales--Jewett arguments is the
following immediate consequence.

\begin{corollary}
  \label{cor:family-lines}
  For every \(k,m\in\N\), with \(k\geq2\), there exists \(N\) such that
  for every family \(\mathcal L\) of combinatorial lines in \(A^N\),
  there is an \(m\)-dimensional combinatorial subspace \(V\subseteq A^N\)
  for which either
  \[
    \operatorname{Lines}(V)\subseteq\mathcal L
  \]
  or
  \[
    \operatorname{Lines}(V)\cap\mathcal L=\varnothing.
  \]
\end{corollary}

\begin{proof}
  Color a line red if it belongs to \(\mathcal L\), and blue otherwise.
  Apply Theorem~\ref{thm:GR-lines} with \(r=2\). The resulting
  \(m\)-dimensional subspace has all of its lines in one color class.
\end{proof}

\end{document}

